\documentclass[tikz,border=8pt]{standalone}
\usepackage{tikz}
\usetikzlibrary{calc}
\usepackage{amsmath}
\usepackage{pgfplots}
\pgfplotsset{compat=1.18}
\usepackage{ctex}

\begin{document}
\begin{tikzpicture}

\begin{axis}[
  width=11cm, height=7cm,
  xmin=-0.02, xmax=1.02,
  ymin=-0.3, ymax=12,
  xlabel={$p$（正面概率）},
  ylabel={后验密度},
  title={\textbf{Beta-Binomial 后验演化（观测 70\% 正面）}},
  axis lines=left,
  xtick={0, 0.2, 0.4, 0.6, 0.7, 0.8, 1.0},
  xticklabel style={font=\scriptsize},
  yticklabel style={font=\scriptsize},
  ytick={0,2,4,6,8,10},
  grid=major, grid style={gray!20},
  clip=false,
  legend pos=north west,
  legend style={font=\small}
]

  % n=0: 先验 Beta(1,1) = Uniform，密度=1
  \addplot[gray!60, thick, domain=0.01:0.99, samples=60] {1.0};
  \addlegendentry{$n=0$（先验 Beta$(1,1)$，均匀）}

  % n=5, 正面3次 -> 后验 Beta(4,3), f ∝ p^3*(1-p)^2
  % 归一化常数 B(4,3) = 3!*2!/6! = 6/720 = 1/120, 1/B = 120/(3*2) = 20... 用数值近似
  % Beta(4,3): f = p^3*(1-p)^2 / B(4,3), B(4,3)=1/60 => f=60*p^3*(1-p)^2
  \addplot[green!60!black, thick, domain=0.01:0.99, samples=80]
    {60*x^3*(1-x)^2};
  \addlegendentry{$n=5$（后验 Beta$(4,3)$）}

  % n=20, 正面14次 -> 后验 Beta(15,7)
  % B(15,7) 数值约 = 1/1144066, 1/B ≈ 1144.066
  % 近似:使用手工归一化，峰值大约在 p=14/20=0.7，
  % f = p^14*(1-p)^6 / B(15,7)
  \addplot[blue!70!black, thick, domain=0.01:0.99, samples=100]
    {5005*x^14*(1-x)^6};
  \addlegendentry{$n=20$（后验 Beta$(15,7)$）}

  % n=100, 正面70次 -> 后验 Beta(71,31)
  % 峰值在 p=70/100=0.7，非常尖锐
  % 使用手工坐标点近似（避免数值溢出问题）
  \addplot[red!70!black, ultra thick] coordinates {
    (0.50, 0.01)(0.55, 0.08)(0.58, 0.28)(0.60, 0.65)(0.62, 1.30)
    (0.64, 2.30)(0.66, 3.60)(0.68, 5.20)(0.70, 6.80)(0.72, 5.20)
    (0.74, 3.40)(0.76, 1.90)(0.78, 0.90)(0.80, 0.38)(0.82, 0.12)
    (0.84, 0.03)(0.86, 0.005)(0.90, 0.0)
  };
  \addlegendentry{$n=100$（后验 Beta$(71,31)$）}

  % 真实 p* 垂线
  \addplot[red!90!black, very thick, dashed] coordinates {(0.7, 0) (0.7, 11.5)};
  \node[font=\small, red!90!black, right] at (axis cs:0.71, 10.5) {$p^* = 0.7$};

  % 注释
  \node[draw, fill=yellow!15, rounded corners=3pt, font=\small, align=center]
    at (axis cs:0.2, 9.5)
    {随 $n$ 增大：\\后验集中于 $p^*=0.7$\\不确定性降低};

\end{axis}

\end{tikzpicture}
\end{document}
